%=============================================================================!
% 7.4  LIOUVILLE'S THEOREM                                                    !
%=============================================================================!
\section{Liouville's theorem}
\label{sec:ha-liouville}

A drop of ink in a shallow, slow stream, whose water moves alike at every
depth, is carried downstream, stretched and folded by the current; the
water cannot be squeezed, and none rises from below or sinks, so the patch
of ink keeps its area while its shape changes. The Hamiltonian
flow does the same to a patch of states. We start many copies of
the pendulum from slightly different states, filling a small region of
the phase plane, a blob, and let each move by Hamilton's equations. This
section shows that the blob keeps its area, whatever the system, and that
drag makes it shrink.

\Cref{fig:ha-blob}(a) follows such a blob of pendulums: an ellipse of
starting states centred on \SI{1.6}{\radian} at rest, reaching
\SI{0.4}{\radian} either side in angle and
\SI{0.8}{\kilogram\metre\squared\per\second} either way in momentum, a
circle of radius 1 stretched by 0.4 along one axis and 0.8 along the
other. Swings of different
size take different times (\cref{sec:os-anharmonic}), so the blob is
sheared into a curved streak; yet its area, $\pi\times0.4\times0.8 = \SI{1.00531}{\joule\second}$,
is the same at every time drawn, to within 2.6 parts in $10^7$, the error
of drawing its edge with 6000 points. An area in this plane is an angle
times an angular momentum, and has the unit of the action of
\cref{sec:la-action}.

\begin{figure}[htb]
\centering
\includegraphics{ch07_hamilton/blob}
\caption{A blob of states of the pendulum of \cref{fig:ha-portrait},
centred on $\theta = \SI{1.6}{\radian}$ at rest, carried by the flow and
drawn at the times given, inside the separatrix (grey). (a) Without drag:
sheared, with its area unchanged. (b) With the drag rate $\gamma =
\SI{0.5}{\per\second}$: shrinking towards the state of rest at the
bottom.}
\label{fig:ha-blob}
\end{figure}

\subsection*{Why the area is kept}

We write the velocity of the flow at the point $(q, p)$ as $(u, w)$, with
$u = \dot{q}$ and $w = \dot{p}$. Each short piece of the edge of the
blob, of length $\Delta s$, moves with the flow, and in a short time
$\tau$ it sweeps out a thin strip. Only the part of the velocity at
right angles to the piece sweeps area, so if $\vb{n}$ is the arrow of
length one at right angles to the piece, pointing out of the blob, the
strip has the area $(u, w)\cdot\vb{n}\,\Delta s\,\tau$, counted
positive when the edge moves outwards. Adding the strips around the whole
edge, the area $\mathcal{A}$ of the blob changes at the rate
\begin{equation*}
   \frac{\dd\mathcal{A}}{\dd t} = \oint (u, w)\cdot\vb{n}\,\dd s ,
\end{equation*}
the flow out through the edge. Going round the edge anticlockwise, a
piece $(\Delta q, \Delta p)$ has the blob on its left, and turning it a
right angle clockwise gives $(\Delta p, -\Delta q)$, which points out of
the blob and has the length $\Delta s$; so $\vb{n}\,\Delta s = (\Delta
p, -\Delta q)$, and
\begin{equation*}
   \frac{\dd\mathcal{A}}{\dd t} = \oint\bigl(u\,\dd p - w\,\dd q\bigr) .
\end{equation*}
Green's theorem (\cref{sec:en-curl}) turns a loop integral into an
integral over the area inside: for any smooth $(F_q, F_p)$, $\oint(F_q\,
\dd q + F_p\,\dd p) = \iint(\partial F_p/\partial q - \partial F_q/\partial
p)\,\dd A$, around the loop anticlockwise. With $F_q = -w$ and $F_p = u$,
\begin{equation}
   \frac{\dd\mathcal{A}}{\dd t} = \iint\Bigl(\frac{\partial u}{\partial q}
   + \frac{\partial w}{\partial p}\Bigr)\dd A .
   \label{eq:ha-area-rate}
\end{equation}
The sum in the large brackets is called the divergence of the flow: by
\cref{eq:ha-area-rate}, applied to a blob so small that the sum hardly
changes across it, it is the rate at which the blob's area grows, divided
by that area. For the Hamiltonian
flow, $u = \partial H/\partial p$ and $w = -\partial H/\partial q$, so
\begin{equation*}
   \frac{\partial u}{\partial q} + \frac{\partial w}{\partial p}
   = \frac{\partial^2H}{\partial q\,\partial p}
   - \frac{\partial^2H}{\partial p\,\partial q} = 0 ,
\end{equation*}
because the two mixed partial derivatives of a smooth function are equal
(\cref{sec:en-curl}). The integrand is zero everywhere, so $\dd\mathcal{A}
/\dd t = 0$: the Hamiltonian flow keeps the area of every region it
carries. This is Liouville's theorem.

With the drag of \cref{sec:os-damped}, $\dot{p} = -\partial H/\partial q -
\gamma p$, so $w$ gains the term $-\gamma p$, and $\partial w/\partial p$
gains $-\gamma$. The divergence is then $-\gamma$ everywhere, and
\cref{eq:ha-area-rate} gives $\dd\mathcal{A}/\dd t = -\gamma\mathcal{A}$,
the equation of \cref{sec:ne-drag}, so
\begin{equation*}
   \mathcal{A}(t) = \mathcal{A}_0\,\mathrm{e}^{-\gamma t} .
\end{equation*}
With $\gamma = \SI{0.5}{\per\second}$, the blob of \cref{fig:ha-blob}(b)
keeps 0.60653, 0.36788 and 0.22313 of its area after 1, 2 and
\SI{3}{\second}, the values of $\mathrm{e}^{-\gamma t}$ to the five
figures printed. Every state ends at rest at the bottom: the drag erases
the differences between starting states, which a Hamiltonian flow never
does.

\subsection*{The stretching of a small square}

Liouville's theorem can also be read off a single state and its close
neighbours, which is the form Chapter~9 uses for methods that move in
steps.

\begin{toolbox}[The Jacobian determinant]
The flow over a time $t$ carries each starting state $(q_0, p_0)$ to the
state $(q, p)$ reached at time $t$, so $q$ and $p$ are functions of $q_0$
and $p_0$. We take a small square of starting states with one corner at
$(q_0, p_0)$ and sides of length $\varepsilon$ along the two axes
(\cref{fig:ha-square}). By the Taylor series to first order, the corner
at $(q_0 + \varepsilon, p_0)$ goes to $(q, p) + \varepsilon\,\vb{a}$, with
$\vb{a} = (\partial q/\partial q_0, \partial p/\partial q_0)$, and the
corner at $(q_0, p_0 + \varepsilon)$ goes to $(q, p) + \varepsilon\,
\vb{b}$, with $\vb{b} = (\partial q/\partial p_0, \partial p/\partial
p_0)$. The far corner $(q_0 + \varepsilon, p_0 +
\varepsilon)$ is a step $\varepsilon$ in both variables, and the
first-order rule $\Delta h \approx (\partial h/\partial x)\,\Delta x +
(\partial h/\partial y)\,\Delta y$ of \cref{sec:en-gradient}, applied to
$q$ and to $p$, takes it to $(q, p) + \varepsilon\vb{a} + \varepsilon
\vb{b}$. To first order, then, the square becomes the parallelogram on the
arrows $\varepsilon\vb{a}$ and $\varepsilon\vb{b}$. The area of the
parallelogram on $\vb{a} = (a_1, a_2)$ and $\vb{b} = (b_1, b_2)$ is the
length of their cross product (\cref{sec:ro-cross}); with third
components zero, the cross product has only the third component $a_1b_2 -
a_2b_1$, so the area is $\lvert a_1b_2 - a_2b_1\rvert$. The four
derivatives form the Jacobian matrix of the flow,
\begin{equation*}
   \vb{J} = \begin{pmatrix}
      \partial q/\partial q_0 & \partial q/\partial p_0\\
      \partial p/\partial q_0 & \partial p/\partial p_0
   \end{pmatrix} ,
\end{equation*}
whose columns are $\vb{a}$ and $\vb{b}$, and the number $a_1b_2 - a_2b_1$
made from a $2\times2$ matrix in this way, the product of the diagonal
entries minus the product of the other two, is its determinant, $\det
\vb{J}$. The square of area $\varepsilon^2$ therefore becomes a
parallelogram of area $\lvert\det\vb{J}\rvert\,\varepsilon^2$: the
determinant is the factor by which the flow stretches small areas. For a
body in free flight, $q = q_0 + p_0t/m$ and $p = p_0$, so
\begin{equation*}
   \vb{J} = \begin{pmatrix} 1 & t/m\\ 0 & 1 \end{pmatrix} ,
   \qquad \det\vb{J} = 1\times1 - \frac{t}{m}\times0 = 1 ,
\end{equation*}
a shear that tilts the square without changing its area
(\cref{ex:ha-shear}). For a map of $n$ variables, the $n\times n$
Jacobian matrix has a determinant too, the factor by which the map
stretches small volumes in $n$ dimensions; its definition for larger
matrices, and this property, are standard results of linear algebra that
we take on trust.
\end{toolbox}

\begin{figure}[htb]
\centering
\begin{tikzpicture}[line cap=round, line join=round, scale=1.15,
   arr/.style={-{Stealth[length=5pt]}, line width=0.8pt}]
   \fill[black!8] (0,0) rectangle (1.6,1.6);
   \draw[line width=0.8pt] (0,0) rectangle (1.6,1.6);
   \draw[arr, accent] (0,0) -- (1.6,0) node[midway, below=2pt]
      {\small $\varepsilon$ along $q_0$};
   \draw[arr, accent] (0,0) -- (0,1.6) node[midway, left]
      {\small $\varepsilon$ along $p_0$};
   \fill (0,0) circle (1.2pt) node[left=3pt] {\small $(q_0, p_0)$};
   \draw[-{Stealth[length=6pt]}, line width=0.6pt, black!50]
      (2.2,0.9) to[bend left=20] node[midway, above] {\small flow} (3.8,0.9);
   \begin{scope}[shift={(4.6,0.0)}]
            \fill[black!8] (0,0) -- (1.8,0.55) -- (2.55,2.2) -- (0.75,1.65) -- cycle;
      \draw[line width=0.8pt] (0,0) -- (1.8,0.55) -- (2.55,2.2) -- (0.75,1.65)
         -- cycle;
      \draw[arr, accent] (0,0) -- (1.8,0.55) node[midway, below right]
         {\small $\varepsilon\vb{a}$};
            \draw[arr, accent] (0,0) -- (0.75,1.65) node[midway, left]
         {\small $\varepsilon\vb{b}$};
      \fill (0,0) circle (1.2pt) node[below left] {\small $(q, p)$};
   \end{scope}
\end{tikzpicture}
\caption{The flow over a time $t$ carries a small square of starting
states into a parallelogram on the arrows $\varepsilon\vb{a}$ and
$\varepsilon\vb{b}$, the columns of the Jacobian matrix times the side
$\varepsilon$; its area is $\lvert\det\vb{J}\rvert$ times that of the
square.}
\label{fig:ha-square}
\end{figure}

Since the Hamiltonian flow keeps the area of every region, a small square
keeps its area $\varepsilon^2$, while the toolbox gives $\lvert\det\vb{J}
\rvert\,\varepsilon^2$ with an error that shrinks faster than
$\varepsilon^2$; letting $\varepsilon$ shrink, $\lvert\det\vb{J}\rvert =
1$. At $t = 0$ the map does nothing and $\det\vb{J} = 1\times1 - 0\times0
= 1$, and $\det\vb{J}$ changes smoothly with $t$, so it cannot jump to
$-1$: $\det\vb{J} = 1$ at every time. With drag, the same reasoning gives
$\det\vb{J} = \mathrm{e}^{-\gamma t}$. Computed with
\texttt{flow\_jacobian}, which follows four copies started a small step
either side of the state along each axis and takes central differences
(\cref{sec:mo-atoms}) of where they arrive, the Jacobian of the pendulum's flow from
$(\theta, p) = (\SI{1}{\radian}, \SI{0.5}{\kilogram\metre\squared\per
\second})$ over \SI{2.5}{\second} has a determinant that differs from 1 by
about
$10^{-10}$, although the map is far from doing nothing: its entries are
1.057 and 0.213 in the first row and 1.157 and 1.179 in the second;
rounded, they give $1.057\times1.179 - 0.213\times1.157 = 0.9998$, which is
1 to the accuracy of the rounding.

For $N$ atoms the same holds in the $6N$ dimensions of phase space. The
divergence of the flow is the sum, over all $6N$ coordinates, of the
derivative of each coordinate's rate with respect to that coordinate,
$\sum_k(\partial\dot{q}_k/\partial q_k + \partial\dot{p}_k/\partial p_k)$;
for Hamilton's equations it is $\sum_k(\partial^2H/\partial q_k\partial
p_k - \partial^2H/\partial p_k\partial q_k)$, zero term by term, and the form of
Green's theorem in many dimensions, the divergence
theorem\cite{Tuckerman2023}, then shows that every region of phase space
keeps its volume and $\det\vb{J} = 1$. A cloud of states can be stretched and folded by
the flow but never compressed. In particular, the number of copies in a
small region that moves with the flow does not change, and neither does
its volume, so the density of copies seen by an observer moving with any
one state stays the same. Since each state also keeps its energy, a cloud
whose density, the number of copies per unit volume, is a function
$f(H)$ of the energy alone does not change at all: the copies now at a
state $\Gamma$ came from a state of the same energy, where the density was
$f\bigl(H(\Gamma)\bigr)$, and have kept that density. Chapter~12 builds statistical mechanics
on such clouds.

\begin{keyidea}
The Hamiltonian flow has zero divergence, because the mixed partial
derivatives of $H$ are equal, so it keeps the area, in many dimensions the
volume, of every region of phase space it carries: the Jacobian
determinant of the flow is 1. Drag at the rate $\gamma$ shrinks areas as
$\mathrm{e}^{-\gamma t}$.
\end{keyidea}

\notebook{07}{watch a blob of states flow, with and without drag}

\begin{selfcheck}
A blob of states of the pendulum with drag at $\gamma =
\SI{0.5}{\per\second}$ has an area of \SI{1}{\joule\second}. What is its
area \SI{4}{\second} later?
\end{selfcheck}
